% ============================================================
%  Chapter 4 — The Robot Brain: Programming and Instructions
%  Robot World: Meet the Machines That Help Us
% ============================================================

\chapter{The Robot Brain}

\chapteropening{%
  A robot can have the strongest arm and the sharpest sensors in the world,
  but without instructions it will just sit there doing nothing.
  The brain is what brings a robot to life.
}
\chapterbadge{Code: The Language Robots Understand}

% --- What is a robot brain? ---
\section*{The Controller}

Inside every robot is a small computer called the \term{controller}.
It reads information from the sensors, decides what to do, and sends
commands to the motors.

\begin{picturethis}
  Imagine if your teacher gave the class no instructions for a science
  experiment. Everyone might mix different things at different times and
  the results would be chaos! A robot programme is like the experiment
  instructions that keep everything on track --- a recipe the robot
  follows step by step.
\end{picturethis}

\character{Dot}{Remember the processor inside a smartphone? A robot's
  controller is very similar \textemdash{} it follows instructions millions
  of times per second.}

\sectionrule

% --- Programming ---
\section*{How Do We Tell Robots What to Do?}

Humans write instructions in a \term{programming language}. Popular ones
for robots include:

\begin{itemize}
  \item \textbf{\term{Python}} \textemdash{} easy to read, great for beginners.
  \item \textbf{C/C++} \textemdash{} fast, used in many industrial robots.
  \item \textbf{\term{Scratch}} \textemdash{} block-based, brilliant for learning logic.
  \item \textbf{\term{ROS}} \textemdash{} Robot Operating System, a toolkit engineers use to
    make different robot parts work together.
\end{itemize}

\begin{funfact}
  The first programmable robot was ``Unimate'', installed in a General Motors
  factory in 1961. It was programmed by manually moving its arm and recording
  the positions \textemdash{} no typing needed!
\end{funfact}

\sectionrule

% --- Algorithms ---
\section*{Algorithms \textemdash{} Step-by-Step Plans}

An \term{algorithm} is a set of clear steps that solves a problem.
Every robot programme is made of algorithms.

\begin{quickmap}
  \textbf{Algorithm: avoid an obstacle}
  \begin{enumerate}
    \item Move forward.
    \item If the distance sensor reads less than 30\,cm, stop.
    \item Turn right 90 degrees.
    \item Go back to step 1.
  \end{enumerate}
\end{quickmap}

\sectionrule

% --- Loops and conditions ---
\section*{Loops and Decisions}

Two superpowers of every programme:

\begin{itemize}
  \item \textbf{\term{Loops}} \textemdash{} repeat something over and over.
    Example: ``Keep spinning the propellers until I say stop.''
  \item \textbf{\term{Conditions}} \textemdash{} choose what to do based on information.
    Example: ``IF the floor sensor sees a cliff, THEN stop and turn around.''
\end{itemize}

\begin{bigidea}
  Every clever thing a robot does is built from simple building blocks:
  sequences, loops, and conditions. Stack enough of them together
  and you get a machine that \emph{looks} intelligent.
\end{bigidea}

\sectionrule

% --- Try It ---
\begin{tryit}
  Write an algorithm for making a sandwich. Be very precise \textemdash{}
  remember, a robot cannot guess! For example:
  \begin{enumerate}
    \item Pick up bread with left gripper.
    \item Place bread on plate.
    \item Pick up butter knife\ldots
  \end{enumerate}
  How many steps did you need?
\end{tryit}

\sectionrule



\begin{experiment}
  \textbf{Programme a human robot.}
  \begin{enumerate}
    \item Ask a friend to be the ``robot'' and blindfold them (or close their eyes).
    \item Place a target (a book or toy) across the room.
    \item You are the programmer --- you may only say: \textbf{Forward},
      \textbf{Stop}, \textbf{Turn left}, or \textbf{Turn right}.
    \item Try to guide your robot to touch the target.
    \item Count how many commands it takes.
  \end{enumerate}
  \textit{Notice how precise you must be --- saying ``go over there''
  means nothing to a robot! This is exactly why programmes need
  clear, step-by-step instructions.}
\end{experiment}

\sectionrule

% --- AI preview ---
\section*{Can Robot Brains Learn?}

Most robots follow fixed programmes. But some newer robots can \emph{learn}
from examples --- a bit like practising until you get better at a video game.

\begin{itemize}
  \item A robot arm might try a thousand ways to pick up an egg until it
    finds the gentlest grip.
  \item A delivery robot might learn which routes have fewer obstacles.
\end{itemize}

This ability is called \term{machine learning}, and it is a type of
\term{artificial intelligence} (AI). We will explore AI much more in
Book~4: \emph{When Machines Start Thinking}!

\sectionrule

% --- Diagram ---
\begin{diagram}[Inside a robot controller: sensors in, decisions made, commands out.]
  \begin{tikzpicture}[>=Stealth, node distance=2.5cm]
    \node[draw=MetalGray,fill=MetalGray!10,rounded corners,
          minimum width=2cm,minimum height=1cm,font=\small\bfseries] (sin) {Sensor data in};
    \node[draw=RoboGreen!80!black,fill=RoboGreen!15,rounded corners=10pt,
          minimum width=3cm,minimum height=1.8cm,line width=1.2pt,
          font=\bfseries,right=of sin] (brain) {Controller};
    \node[draw=WarmOrange,fill=WarmOrange!10,rounded corners,
          minimum width=2cm,minimum height=1cm,font=\small\bfseries,right=of brain] (mout) {Motor commands};
    \draw[->,line width=1.5pt,MetalGray] (sin) -- (brain);
    \draw[->,line width=1.5pt,WarmOrange] (brain) -- (mout);
    \node[font=\scriptsize\itshape,below=0.3cm of brain] {Runs the programme};
  \end{tikzpicture}
\end{diagram}

\sectionrule


\begin{quickcheck}
  \begin{enumerate}
    \item What is an algorithm?
    \item Name the two programming superpowers mentioned in this chapter.
    \item Why does a robot need very precise instructions?
  \end{enumerate}
\end{quickcheck}

\sectionrule
% --- New Words ---
\begin{newwords}
  \begin{description}[labelwidth=3.8cm, leftmargin=4.0cm]
    \item[Controller] The computer board inside a robot that runs the programme.
    \item[Programme] A set of instructions written in a programming language.
    \item[Algorithm] A step-by-step method for solving a problem.
    \item[Loop] An instruction that repeats until a condition is met.
    \item[Condition] A yes/no check that changes what the programme does next.
    \item[ROS] Robot Operating System \textemdash{} a popular framework for robot software.
  \end{description}
\end{newwords}
